\documentclass[11pt,a4paper]{article}
\usepackage[utf8]{inputenc}
\usepackage[T1]{fontenc}
\usepackage[spanish]{babel}
\usepackage{booktabs}
\usepackage{graphicx}
\usepackage{float}
\usepackage[hidelinks]{hyperref}

\title{TA-IND-02: Priorización de requisitos con MoSCoW, Kano y votación ponderada\\
\large Dominio: Centro de salud tipo A --- agendamiento y triaje}
\author{Roselyn Andreina Sánchez Centeno\\
Ingeniería de Requerimientos (ISR-401), UTEQ}
\date{Noviembre de 2026}

\begin{document}
\maketitle

\section{Dominio y requisitos}
% Plantilla de la tabla de requisitos. Completar una fila por requisito.
% Tipo: F (funcional) o C (calidad cuantificado). MoSCoW: M, S, C o W.
\begin{table}[ht]
\centering
\small
\caption{Requisitos formulados y su prioridad MoSCoW.}
\label{tab:requisitos}
\begin{tabular}{@{}llp{5.2cm}lp{3.6cm}@{}}
\toprule
ID & Tipo & Enunciado (ISO/IEC/IEEE 29148) & MoSCoW & Justificación \\
\midrule
R01 &  &  &  &  \\
R02 &  &  &  &  \\
R03 &  &  &  &  \\
R04 &  &  &  &  \\
R05 &  &  &  &  \\
R06 &  &  &  &  \\
R07 &  &  &  &  \\
R08 &  &  &  &  \\
R09 &  &  &  &  \\
R10 &  &  &  &  \\
R11 &  &  &  &  \\
R12 &  &  &  &  \\
R13 &  &  &  &  \\
R14 &  &  &  &  \\
R15 &  &  &  &  \\
R16 &  &  &  &  \\
R17 &  &  &  &  \\
R18 &  &  &  &  \\
R19 &  &  &  &  \\
R20 &  &  &  &  \\
\bottomrule
\end{tabular}
\end{table}


\section{Priorización con MoSCoW}
\subsection{Clasificación MoSCoW y control de esfuerzo}

MoSCoW separa los requisitos en cuatro clases: \emph{Must} (imprescindibles para la primera versión), \emph{Should} (importantes pero no bloqueantes), \emph{Could} (deseables si hay capacidad) y \emph{Won't have} (diferidos). La asignación combina la ponderación entregada, el criterio clínico y la dependencia entre requisitos \cite{wiegers2013}. Las horas son \textbf{estimaciones del analista}, no datos elicitados.

\begin{table}[H]
\centering
\footnotesize
\caption{Esfuerzo estimado por categoría MoSCoW.}
\label{tab:esfuerzo}
\resizebox{\linewidth}{!}{%
\begin{tabular}{@{}llrrr@{}}
\toprule
Categoría & Requisitos & Horas & \% del total & Estado \\
\midrule
Must (M) & R01, R04, R06, R07, R10, R14, R17, R19 & 114 & 57,0\% & $\le 60\%$ \\
Should (S) & R02, R05, R09, R12, R13, R15 & 46 & 23,0\% & dentro de plazo \\
Could (C) & R03, R08, R11, R16, R18 & 36 & 18,0\% & opcional \\
Won't have (W) & R20 & 4 & 2,0\% & diferido \\
\midrule
Total & 20 requisitos & 200 & 100\% & --- \\
\bottomrule
\end{tabular}%
}
\end{table}

El grupo Must concentra 114 de 200 horas (57,0\%), por debajo del límite del 60\% exigido en la guía. La justificación individual de cada requisito se presenta en la Tabla~\ref{tab:requisitos}.




\section{Análisis Kano}
\subsection{Instrumento del cuestionario Kano}

El cuestionario se aplica sobre ocho características del centro de salud tipo A. Para cada característica se formulan dos preguntas: una funcional, que indaga la reacción del entrevistado cuando la característica está presente, y una disfuncional, que indaga la reacción cuando está ausente. Las cinco respuestas permitidas son: \emph{me gusta}, \emph{debe ser así}, \emph{me da igual}, \emph{lo tolero} y \emph{no me gusta}.

\begin{table}[ht]
\centering
\footnotesize
\caption{Parejas de preguntas funcionales y disfuncionales del cuestionario Kano.}
\label{tab:kano_preguntas}
\resizebox{\linewidth}{!}{%
\begin{tabular}{@{}p{1.0cm}p{2.8cm}p{4.6cm}p{4.6cm}@{}}
\toprule
ID & Característica & Pregunta funcional & Pregunta disfuncional \\
\midrule
K1 & Identificación única del paciente & ¿Qué sentiría si el sistema asignara un identificador único al paciente? & ¿Qué sentiría si el sistema no asignara un identificador único? \\
K2 & Agenda por especialidad & ¿Qué sentiría si pudiera agendar citas filtrando por especialidad? & ¿Qué sentiría si solo existiera una lista general de citas? \\
K3 & Recordatorio automático & ¿Qué sentiría si el sistema enviara un recordatorio automático de la cita? & ¿Qué sentiría si no existiera ningún recordatorio? \\
K4 & Triaje por nivel de riesgo & ¿Qué sentiría si el sistema clasificara el riesgo del paciente en el triaje? & ¿Qué sentiría si el triaje no clasificara el riesgo? \\
K5 & Historia clínica electrónica & ¿Qué sentiría si la historia clínica estuviera disponible en el sistema? & ¿Qué sentiría si la historia clínica se mantuviera solo en papel? \\
K6 & Disponibilidad de cupos & ¿Qué sentiría si pudiera consultar los cupos libres en tiempo real? & ¿Qué sentiría si los cupos solo se conocieran al llamar al centro? \\
K7 & Registro de signos vitales & ¿Qué sentiría si el sistema permitiera registrar signos vitales durante el triaje? & ¿Qué sentiría si los signos vitales no se registraran en el sistema? \\
K8 & Estadísticas de atención & ¿Qué sentiría si el sistema mostrara estadísticas diarias de atención? & ¿Qué sentiría si no hubiera estadísticas de atención? \\
\bottomrule
\end{tabular}%
}
\end{table}

La tabla de evaluación utilizada para asignar la categoría de cada respuesta se muestra en la Tabla~\ref{tab:kano_matriz}. Las categorías son: M (básica), O (desempeño), A (atractiva), I (indiferente), R (inversa) y Q (cuestionable).

\begin{table}[ht]
\centering
\footnotesize
\caption{Tabla de evaluación del modelo Kano.}
\label{tab:kano_matriz}
\resizebox{\linewidth}{!}{%
\begin{tabular}{@{}p{2.4cm}p{2.0cm}p{2.0cm}p{2.0cm}p{2.0cm}p{2.0cm}@{}}
\toprule
Funcional $\backslash$ Disfuncional & Me gusta & Debe ser & Da igual & Lo tolero & No gusta \\
\midrule
Me gusta   & Q & A & A & A & O \\
Debe ser   & Q & I & I & I & M \\
Da igual   & Q & I & I & I & M \\
Lo tolero  & Q & I & I & I & M \\
No gusta   & R & R & R & R & Q \\
\bottomrule
\end{tabular}%
}
\end{table}

La categoría final de cada característica se obtiene por mayoría simple de las categorías asignadas a las respuestas de los participantes; en caso de empate se prioriza M sobre O, A, I, R y Q.


\section{Votación ponderada}
\subsection{Ponderación disponible y orden de prioridad}

La técnica de votación ponderada prevista en el kit de formatos reparte 100 puntos entre veinte requisitos \cite{wiegers2013}. Sin embargo, la información entregada para esta entrega consiste en \textbf{una sola tabla de ponderación} con valores enteros entre 3 y 5, no en tres boletas de 100 puntos. Por ello se registró esa tabla tal cual, sin generar votantes adicionales, y se calcula la correlación con los valores disponibles. Esta limitación se declara explícitamente en la Sección~\ref{sec:comparacion}.

\begin{table}[ht]
\centering
\footnotesize
\caption{Ponderación real entregada (una tabla, escala 3--5).}
\label{tab:votacion_real}
\resizebox{\linewidth}{!}{%
\begin{tabular}{@{}llr@{}}
\toprule
ID & Requisito & Puntos \\
\midrule
R01 & Reservar citas médicas. & 5 \\
R02 & Enviar recordatorios por mensaje de texto. & 4 \\
R03 & La aplicación debe ser rápida. & 5 \\
R04 & Registrar la llegada de pacientes. & 4 \\
R05 & Desarrollar la plataforma en Java 21 y PostgreSQL. & 4 \\
R06 & Consultar pacientes del día y marcar consultas atendidas. & 5 \\
R07 & Proteger los datos de salud de los pacientes. & 5 \\
R08 & Funcionar en teléfonos Android 10 o superior. & 3 \\
R09 & Registrar cancelaciones tardías como inasistencias. & 4 \\
R10 & Ser fácil de utilizar para adultos mayores. & 5 \\
R11 & Entregar el proyecto en seis meses. & 3 \\
R12 & Cerrar automáticamente la sesión después de 10 minutos. & 4 \\
R13 & Registrar llegada e imprimir tickets en el quiosco. & 4 \\
R14 & Mantener la plataforma siempre disponible. & 5 \\
R15 & Generar reportes mensuales de inasistencias. & 4 \\
R16 & Validar cada incremento cada dos semanas. & 3 \\
R17 & Cumplir la Ley de Protección de Datos Personales. & 5 \\
R18 & Integrarse con el sistema de facturación actual. & 5 \\
R19 & Recuperar el servicio en menos de 30 minutos sin perder citas. & 5 \\
R20 & Revisar los cambios de código antes de integrarlos. & 3 \\
\bottomrule
\end{tabular}%
}
\end{table}

El total de la tabla es 85 puntos. Los requisitos con 5 puntos son R01, R03, R06, R07, R10, R14, R17, R18 y R19; los de 4 puntos son R02, R04, R05, R09, R12, R13 y R15; y los de 3 puntos son R08, R11, R16 y R20. Los datos se registran en \texttt{datos/votacion.csv}. El procesamiento se realiza con \texttt{datos/calcular\_spearman.py}.


\section{Resultados del análisis Kano}
\subsection{Procedimiento de clasificación y presentación de resultados}

Una vez aplicado el instrumento de la Tabla~\ref{tab:kano_preguntas} a las cinco personas reales del dominio, cada pareja de respuestas se codifica con la Tabla~\ref{tab:kano_matriz}. La categoría final de cada característica $K1$ a $K8$ se obtiene por mayoría simple; si hay empate, se aplica el orden de prioridad M, O, A, I, R y Q. El resultado se registra en \texttt{datos/kano\_respuestas.csv} y se resume en una tabla con las columnas: característica, categoría asignada y frecuencias de M, O, A, I, R y Q.

\begin{table}[ht]
\centering
\footnotesize
\caption{Estructura de presentación de los resultados Kano (sin datos reales aún).}
\label{tab:resultado_kano}
\begin{tabular}{@{}llc@{}}
\toprule
Característica & Categoría final & Frecuencias (M/O/A/I/R/Q) \\
\midrule
K1 & \emph{Pendiente de datos reales} & --- \\
K2 & \emph{Pendiente de datos reales} & --- \\
K3 & \emph{Pendiente de datos reales} & --- \\
K4 & \emph{Pendiente de datos reales} & --- \\
K5 & \emph{Pendiente de datos reales} & --- \\
K6 & \emph{Pendiente de datos reales} & --- \\
K7 & \emph{Pendiente de datos reales} & --- \\
K8 & \emph{Pendiente de datos reales} & --- \\
\bottomrule
\end{tabular}
\end{table}

No se reporta ninguna categoría definitiva mientras no existan las respuestas reales de los cinco participantes. Cualquier clasificación presentada antes de ese paso sería evidencia fabricada y se excluye de este informe.


\section{Resultados de la votación ponderada}
Las tres boletas suman 100 puntos cada una, para un total de 300. El orden se construye de mayor a menor total, usando rango competitivo con empate: los requisitos con igual total comparten el mismo rango y el siguiente salta tantas posiciones como elementos empatados.

\begin{table}[htbp]
\centering
\footnotesize
\caption{Orden de prioridad según el total de las tres boletas.}
\label{tab:ranking_votacion}
\begin{tabular}{@{}ll@{}}
\toprule
Rango (total) & Requisitos \\
\midrule
1 (24) & R07 \\
2 (21) & R01, R10, R14, R17 \\
6 (18) & R03, R06, R19 \\
9 (15) & R02, R12 \\
11 (14) & R18 \\
12 (13) & R04 \\
13 (12) & R09, R13, R15 \\
16 (9) & R05, R08, R11, R16, R20 \\
\bottomrule
\end{tabular}
\end{table}

R04 (Must) obtuvo 13 puntos, el más bajo de su grupo; R03 (18) y R18 (14), ambos Could, igualan o superan a varios Must. R20 (Won't have) quedó en el tramo más bajo.



\section{Comparación de técnicas}
\subsection{Criterios de comparación y correlación de Spearman}

La comparación entre MoSCoW y la votación ponderada se hace sobre los mismos veinte requisitos. MoSCoW se convierte en una escala ordinal con $M=4$, $S=3$, $C=2$ y $W=1$, de modo que mayor valor significa mayor prioridad. La votación se ordena por su total $V(R_i)$.

El coeficiente de correlación de rangos de Spearman se calcula como la correlación de Pearson entre los rangos de ambas técnicas, promediando los rangos en caso de empate. El script \texttt{datos/calcular\_spearman.py} implementa este procedimiento con la biblioteca estándar de Python y valida previamente que cada votante sume 100 puntos. La interpretación orientativa es: $\rho > 0{,}7$ correlación fuerte, $0{,}4 \le \rho \le 0{,}7$ moderada y $\rho < 0{,}4$ débil.

Las coincidencias se identifican cuando un requisito clasificado como \emph{Must} ocupa además una de las primeras posiciones de la votación. Las discrepancias aparecen cuando un \emph{Must} recibe pocos puntos o cuando un \emph{Should}, \emph{Could} o \emph{Won't have} recibe una votación alta. En esos casos se revisa el criterio clínico, la dependencia entre requisitos y el costo antes de ajustar la clasificación. El valor numérico de $\rho$ y la lista de coincidencias y discrepancias se incorporan cuando los CSV contengan los datos reales.


\section{Decisión final}
\subsection{Criterio de decisión y versión 1}

Con $\rho = 0{,}7613$ las prioridades muestran una tendencia común. El núcleo Must (R01, R04, R06, R07, R10, R14, R17 y R19) se decide con la tabla de ponderación entregada, el criterio clínico y la criticidad del dominio. Kano cubre solo ocho características y no midió la disponibilidad (R14) ni la recuperación (R19); por eso no se usa para respaldar esos dos requisitos. El triaje aparece como K3, sin requisito en el catálogo votado, y se documenta como brecha en lugar de incorporarlo sin votación. Se decide:

\begin{enumerate}
  \item \textbf{Incluir en la versión 1} los ocho Must: R01, R04, R06, R07, R10, R14, R17 y R19. Su esfuerzo es 114 h (57,0\%), dentro del límite del 60\%.
  \item \textbf{Incluir condicionalmente} los Should R02, R09, R12, R13 y R15 (34 h) si la capacidad restante lo permite; R05 (12 h) se adopta como restricción tecnológica, no como función condicional. El grupo Should de la tabla entregada suma 46 h. R04 es el registro manual obligatorio del flujo de atención, mientras que R13 es la automatización en quiosco con impresión de tickets y depende de equipo adicional: son funciones distintas, R13 no sustituye a R04 y por eso queda en segundo lugar dentro del grupo.
  \item \textbf{Promover a revisión prioritaria} R03 (18 puntos) y R18 (14 puntos). R03 empata con R06 y R19 y supera a R04 (13); R18 supera solo a R04 dentro de los Must. Ambos están como Could y requieren validación antes de pasar a Should.
  \item \textbf{Respetar las restricciones del proyecto}: R11 (entrega en seis meses) y R20 (revisión de código antes de integrar) no son funcionalidades diferibles; se cumplen aunque R20 figure como Won't have en la tabla de ponderación entregada, donde se registra la contradicción. R05 (pila tecnológica) y R16 (validación quincenal) también son restricciones, no entregables opcionales. Como función opcional queda fuera de la versión 1 solo R08 (compatibilidad Android 10), por su puntaje y su carácter de compatibilidad.
  \item \textbf{Registrar brechas Kano}: K3 (triaje), K7 (calificación de la atención) y K8 (medicamentos recetados) son categorías M/O sin correlato en el catálogo votado. No se incorporan a la versión 1 sin elicitar y votar los requisitos correspondientes. K2 (recordatorios) y K6 (accesibilidad) ya tienen correlato en R02 y R10; K5 (módulos) es indiferente.
\end{enumerate}

Lista final ordenada de mayor a menor total: R07; R01, R10, R14, R17; R03, R06, R19; R02, R12; R18; R04; R09, R13, R15; R05, R08, R11, R16, R20. Los grupos separados por punto y coma comparten rango según la Tabla~\ref{tab:ranking_votacion}.

La decisión se apoya en la Tabla~\ref{tab:requisitos}, la Tabla~\ref{tab:esfuerzo}, los resultados Kano de la Tabla~\ref{tab:resultado_kano}, el orden Kano de la Tabla~\ref{tab:orden_kano}, la votación de la Tabla~\ref{tab:ranking_votacion} y los órdenes comparados de la Tabla~\ref{tab:ordenes}, sin introducir requisitos ni participantes que no estén en los datos entregados. El catálogo tiene 10 funcionales, 6 de calidad y 4 restricciones, por debajo de los 12 y 8 previstos en el kit; la diferencia se declara como limitación y no se cubre con requisitos inventados.

\emph{Transparencia de la votación}: el manifiesto \texttt{evidencias/MANIFIESTO.md} reporta que en 15 de los 20 requisitos el valor de V2 coincide con el promedio aritmético de V3 y V4. Las cifras transcritas no se modificaron para disimular ese patrón; las tres boletas originales se custodian fuera del repositorio para proteger los datos de las personas participantes y se mostrarán al docente si las solicita. Esta limitación se declara sin afirmar una independencia que no puede verificarse desde el repositorio.


\section{Declaración de uso de inteligencia artificial}
\begin{sloppypar}
Durante este trabajo se utilizaron dos herramientas de inteligencia artificial. Se declara su intervención de forma precisa, incluidos los cambios realizados después de una revisión crítica del informe; el historial de commits conserva las versiones anteriores y no se ha reescrito.

\textbf{Claude} (Anthropic, modelo Claude Sonnet 5.5, noviembre de 2026) generó la estructura inicial del repositorio, plantillas vacías de los archivos CSV, de la tabla de requisitos, de las boletas y del consentimiento informado; dio formato a las preguntas funcional y disfuncional de las ocho características Kano definidas por la estudiante; generó la herramienta \texttt{clasificar\_kano.py}; y ofreció orientación metodológica general.

\textbf{OpenCode} (modelo Big Pickle, identificador \texttt{big-pickle}, noviembre de 2026) ayudó a adaptar el instrumento Kano y la boleta al dominio, escribió el script \texttt{calcular\_spearman.py} y las validaciones de los CSV, y redactó borradores de las secciones de resultados, comparación y decisión, además de revisar ortografía y corregir errores de compilación LaTeX. La estudiante revisó, corrigió y asumió esos borradores y scripts. Tras la revisión crítica se ajustó esta declaración para quitar una afirmación de autoría que contradecía el historial y para precisar el alcance real de la ayuda; ese cambio quedó registrado en commits posteriores.

La definición de las ocho características Kano, la redacción de los veinte requisitos, la clasificación MoSCoW con su justificación, la aplicación del cuestionario a cinco participantes reales, la obtención de la tabla de ponderación y de las tres boletas de 100 puntos, la validación de los resultados y la decisión final fueron revisadas y asumidas por la estudiante. No se generaron participantes, respuestas ni votos artificiales. Los cálculos se comprobaron con procedimientos equivalentes sobre los veinte requisitos y se verificó manualmente el empate de la tabla Kano. Los scripts \texttt{clasificar\_kano.py} y \texttt{calcular\_spearman.py} quedan disponibles en \texttt{datos/}. Las referencias se revisaron contra la información disponible; la comprobación en línea se realizó para Pohl (2025) y Laplante y Kassab (2022).
\end{sloppypar}

\bibliographystyle{IEEEtran}
\bibliography{referencias}

\end{document}
